\section{Related Work}
\label{sec:related}

\para{Curriculum learning for language models.}
\citet{bengio2009curriculum} introduced curriculum learning---training on progressively harder examples---and demonstrated improved convergence on language modeling tasks. Subsequent work applied curricula to language model pretraining with mixed results. \citet{campos2021corpus} tested eight difficulty-based curricula for language modeling and found no compelling benefit, concluding that naive difficulty orderings do not help at scale. More recently, \citet{fan2023irreducible} proposed an \emph{irreducible curriculum} that uses a proxy model to score sample learnability, achieving consistent perplexity improvements across all domains of RedPajama. \citet{chen2023skillit} introduced a skills-based framework that models prerequisite dependencies between data subsets, achieving 3$\times$ data efficiency on RedPajama. A critical methodological insight comes from \citet{luo2025lrdecay}, who showed that standard learning rate decay schedules mask curriculum benefits by reducing the learning rate precisely when the best data is presented; they recommend constant learning rate with model averaging, which we adopt. Unlike these approaches, which focus on difficulty-based or learnability-based sample ordering, we study \emph{domain-level} distribution shifts where entire domains are withheld and later introduced.

\para{Data mixture optimization.}
Rather than ordering data temporally, another line of work optimizes the static mixture of training domains. \citet{xie2023doremi} proposed DoReMi, which uses a small proxy model trained with Group DRO to find optimal domain weights, achieving 2.6$\times$ faster convergence on downstream tasks. \citet{lin2024rho1} introduced token-level selection via excess loss, training only on the most informative tokens. These approaches optimize \emph{what} to train on but not \emph{when}---they do not consider the temporal ordering of domains. Our work is complementary: \dsc could be combined with optimized mixture weights within each phase.

\para{Distribution shift robustness.}
A separate literature studies how to train models robust to distribution shifts at test time. \citet{arjovsky2019irm} proposed Invariant Risk Minimization, learning representations where the optimal classifier is invariant across environments. \citet{sagawa2019dro} showed that Group DRO with strong regularization improves worst-group performance under distribution shift. \citet{krueger2021rex} introduced Risk Extrapolation, penalizing variance of per-domain risks. \citet{yi2022pretrain} found that pretraining on diverse data provides substantial distribution shift robustness, often more than specialized methods applied during fine-tuning. Our work connects these two literatures: rather than applying robustness methods at fine-tuning time, we design the pretraining curriculum itself to expose the model to controlled distribution shifts, encouraging rapid adaptation as an emergent property.

\para{Meta-learning and fast adaptation.}
The idea of training for fast adaptation originates in meta-learning. MAML~\cite{finn2017maml} trains models to reach good performance on new tasks within a few gradient steps. While meta-learning typically operates over explicit task distributions with outer-loop optimization, our approach achieves a similar effect implicitly: the Phase~1$\to$Phase~2 transition is itself a distribution shift that the model ``learns from'' during standard gradient descent. This connection to meta-learning is conceptual rather than mechanistic---we do not use bi-level optimization---but the empirical result (faster adaptation) aligns with meta-learning objectives.
